\documentclass[aspectratio=169,usepdftitle=true,11pt,ngerman,t]{beamer}

\newcommand\lqfirstname{Name}
\newcommand\lqfullname{Name Last}
\newcommand\lqemail{name.last}
\errorcontextlines999999

% base
\usepackage[utf8]{inputenc}
\usepackage[T1]{fontenc}
\usepackage{lmodern}
\usepackage{microtype}
\usepackage[main=ngerman]{babel}
\usepackage{csquotes}

% graphics
\usepackage[prefix=]{xcolor-material}
\usepackage{graphics}
\usepackage{tcolorbox}
\usepackage{twemojis}
\usepackage{tikz}
\usetikzlibrary{
    positioning,
    calc,
    tikzmark,
    circuits.logic.IEC
}

% math
\usepackage{amsmath,amssymb}
\usepackage{mathtools}
\usepackage{array}
\usepackage{booktabs}
\usepackage{cancel}
\usepackage{ulem}

% flo
\usepackage{fancyqr}
\usepackage[addons]{color-palettes}
\usepackage[encoding,noloadlangs,cpalette,numinpar,fakeminted]{sopra-listings}
\usepackage{code-animation}
\usetheme[theorems]{colorful-dream}
\theoremstyle{exercise}
\newtheorem*{exercise*}{Aufgabe}
\theoremstyle{solve}
\newtheorem*{solve*}{Lösung}

% tools
\usepackage{algpseudocode}
\usepackage{dirtree}

% warnings
\usepackage{silence}
\WarningFilter{latexfont}{Some font shapes were not available}
\WarningFilter{latexfont}{Font shape}
\WarningFilter{pgf}{Returning node center instead of a point on node border}
\let\origenddocument\enddocument\renewcommand{\enddocument}{\hfuzz=10.95pt\origenddocument}
% title image: box
% ([vshift], inverse scale, width, content)
\newcommand\titleimagebox[4][0pt]{%
\newlength\tibheight\setlength\tibheight{4cm}
\newlength\tibtY\setlength\tibtY{-6em}
\titleimage{\begin{minipage}[c][#2\tibheight]{#3}\vspace*{#2\tibtY}%
    \vspace*{#1}\begin{tcolorbox}[boxrule=0.1pt]%
    \centering#4% TODO: bug: not centered if too wide
    \end{tcolorbox}%
\end{minipage}}}

% outro image: text
% (text)
\newcommand\outroimagetext[1]{\outro{Karlsruhe, \today%
\begin{tikzpicture}[remember picture,overlay]
    \node[align=center] at (current page.center) {#1};
\end{tikzpicture}}}% chktex 31

% reset counters
\newcommand\resetframecounters{\addtocounter{exercise}{-1}\addtocounter{solve}{-1}}
\newcommand\resetex{\addtocounter{exercise}{-1}}
\newcommand\resetsolve{\addtocounter{solve}{-1}}

% small footnote
\newcommand{\nicefootnote}[1]{\tikz[overlay,remember picture]{\node[anchor=south,text width=\dimexpr\textwidth+2em\relax,align=center,above=1em of current page.south] (nicefootnote) {\hrule\vspace{1pt}\footnotesize#1};}}

% kv
% ([draw], variables, values)
% {currently only works correctly for 4 variables}
\newcommand{\lqkv}[3][]{
    \begin{tikzpicture}
        \foreach \v [count=\k from 1] in {#2} {\global\let\Vdim=\k}
        \pgfmathsetmacro\rows{2^ceil(\Vdim/2)}
        \pgfmathsetmacro\cols{2^int(\Vdim/2)}

        \foreach \v [count=\k from 1] in {#2} {
            \pgfmathsetmacro\shift{0.5+floor(\k*ceil(\Vdim/2)/2)}

            \ifcase\k\relax
            \or
                \draw (\shift-1,0.75) -- (\shift+1,0.75);
                \node[fill=white] at (\shift,0.75) {\v};
            \or
                \draw (-0.75,-\shift) -- (-0.75,-\shift+2);
                \node[fill=white] at (-0.75,-\shift+1) {\v};
            \or
                \draw (\shift-2,-\cols+0.25) -- (\shift,-\cols+0.25);
                \node[fill=white] at (\shift-1,-\cols+0.25) {\v};
            \or
                \draw (\cols-0.25,-\shift+1) -- (\cols-0.25,-\shift+3);
                \node[fill=white] at (\cols-0.25,-\shift+2) {\v};
            \fi
        }
      
        \foreach \v [count=\k from 0] in {#3} {
            \pgfmathsetmacro\x{mod(\k,\rows)}
            \pgfmathsetmacro\y{int(\k/\rows)}
            \node[draw, minimum width=1cm, minimum height=1cm, align=center] at (\x,-\y) {\v};
        }

        \begin{scope}
            \clip (-0.5,0.5) rectangle (3.5,-3.5);
            \foreach \r in {#1} {\draw[rounded corners=1ex] \r;}
        \end{scope}

        \global\let\Vdim\undefined
    \end{tikzpicture}
}
\title{Digitaltechnik}
\institute{Karlsruher Institut für Technologie}

\author{\lqfullname}
\email{\lqemail @student.kit.edu}% chktex 1
% credit flo
\addtobeamertemplate{title page}{}{%
\scriptsize\color{gray}\begin{tikzpicture}[overlay, remember picture]
    \node[above left = 1em and 1em of current page.south east,align=center] (n0) {\LaTeX-Vorlage von\\\href{https://github.com/EagleoutIce}{Florian Sihler}};
    \node[above = 0cm of n0] (n1) {\fancyqr[size=1.5cm,tight]{https://eagleoutice.github.io/beamer-themes/\#colorful-dream}};
\end{tikzpicture}}% chktex 31
\usepackage{tikz}
\usetikzlibrary{fadings}

% cursed braces
\makeatletter
\newcommand{\cb@dx}[1]{%
    \ifnum`#1=`(  .0652em\else% chktex 37
    \ifnum`#1=`) -.0652em\else
    \ifnum`#1=`<  .0652em\else
    \ifnum`#1=`> -.0652em\else
    \ifnum`#1=`B  .0113em\else
    0em\fi\fi\fi\fi\fi
}
\newcommand{\cb@r}[1]{%
    \ifnum`#1=`< \langle\else
    \ifnum`#1=`> \rangle\else
    \ifnum`#1=`& \backslash\else
    \ifnum`#1=`= \forall\else% im not sure about this
    #1\fi\fi\fi\fi
}
\newcommand{\cb@clip}[4][]{\scoped{%
    \clip(-.3117em,#2) rectangle (.3117em,#3);%
    \node[#1] at (\cb@dx{#4},0em) {\(\cb@r{#4}\)};% chktex 1
}}
\newcommand{\cb@ht}[1][]{\cb@clip[#1]{0em}{.5669em}}
\newcommand{\cb@hb}[1][]{\cb@clip[#1]{-.5669em}{0em}}
\newcommand{\cb@tt}[1][]{\cb@clip[#1]{.1418em}{.5669em}}
\newcommand{\cb@tm}[1][]{\cb@clip[#1]{-.1418em}{.1418em}}
\newcommand{\cb@tb}[1][]{\cb@clip[#1]{-.5669em}{-.1418em}}
\newcommand{\cb@split}[4][]{%
    \ifx\relax#4\relax
        \cb@ht[#1]{#2}\cb@hb[#1]{#3}%
    \else
        \cb@tt[#1]{#2}\cb@tm[#1]{#3}\cb@tb[#1]{#4}%
    \fi
}
\newcommand{\cb}[2][]{{\tikz[baseline=-.8ex]{\cb@split[#1]#2\relax}}}
\newcommand{\cblogo}{\cb[inner sep=0pt,shading=axis,left color=orange,right color=violet,shading angle=0]{cB}}
\makeatother

\usepackage{circuitikz}% chktex-file 1 chktex-file 8 chktex-file 26
\usetikzlibrary{overlay-beamer-styles}

%%%%%%%%%%%%%%%%%%%%%%%%%%%%%%%%%%%%%%%%%%%%%%%%%%%%%%%%%%%%%%%%

\subtitle{Tutorium 6}
\date{23. Juni 2026}

\titleimagebox{0.5}{2cm}{\texttwemoji{sun}\texttwemoji{thermometer}\texttwemoji{hot_face}}

\outroimagetext{I wish I had Hitzefrei!}

%%%%%%%%%%%%%%%%%%%%%%%%%%%%%%%%%%%%%%%%%%%%%%%%%%%%%%%%%%%%%%%%

\begin{document}% TODO: Animation

\section[Korrektur des Blatts, das ich noch nicht korrigiert habe:]{Korrektur}

\subsection{Aufgabe 1}
\begin{frame}{Diesmal habe ich die Zeichnung!}\onslide<+->
    \onslide<+->
    \vspace*{-1.5\abovedisplayskip}
    \begin{columns}[T]
    \begin{column}{0.5\textwidth}
        \begin{exercise}
            Was ist falsch?
            \only<4->{\\Warum?}
            \only<5->{\\Welche Funktion wird dargestellt?}
        \end{exercise}
        \only<6->{
            \begin{solve}
                \centering \(g(d,c,b,a) = \cb{(\{}a \land d\cb{\}]} \lor \cb{<\}[}a \land b\cb{&]} \lor c\)% chktex 9
            \end{solve}
            \resetsolve%
        }
    \end{column}
    \begin{column}{0.5\textwidth}
        \only<-4>{
        \vspace*{\abovedisplayskip}
        \resizebox{\textwidth}{!}{
            \begin{circuitikz}[european]
                % Versorgungsspannung
                \draw[alt=<3->{red}{black}] (0,0) node[ocirc] {} node[anchor=east] {\(V_{dd}\)} -- ++ (0,-1) node[circ] (dot_vdd) {};
                
                %linke Seite
                \draw (dot_vdd) -- ++ (-4,0) -- ++ (0,-0.5) node[circ] (p_dot_ab) {};
                \draw (p_dot_ab) -- ++ (-1,0) node[pfet, anchor=drain, yscale=-1] (pmos_a) {};
                \draw (p_dot_ab) -- ++ (1,0) node[pfet, anchor=drain, yscale=-1] (pmos_b) {};
                \draw (pmos_a.source) -- ++ (1,0) node[circ] (p_dot_ab_2) {};
                \draw (pmos_b.source) -- (p_dot_ab_2);
                
                \draw (p_dot_ab_2) -- ++ (0,0) node[pfet, anchor=drain, yscale=-1] (pmos_c) {};
                \draw (pmos_c.source) -- ++ (0,0) node[circ] (cmos_1) {};
                
                \draw[alt=<3->{red}{black}] (cmos_1) -- ++ (0,-0.5) node[circ] (n_dot_ab) {};
                \draw (n_dot_ab) -- ++ (-1,0) node[nfet, anchor=drain] (nmos_a) {};
                \draw (nmos_a.source) -- ++ (0,0) node[nfet, anchor=drain] (nmos_b) {};
                
                \draw[alt=<3->{red}{black}] (n_dot_ab) -- ++ (1,0) -- ++ (0,-0.75) node[nfet, anchor=drain] (nmos_c) {};
                \draw[alt=<3->{red}{black}] (nmos_c.source) -- ++ (0,-0.78) node[circ] (n_dot_c) {};
                \draw (nmos_b.source) -- (n_dot_c)  {};
                
                % Mitte
                
                \draw[alt=<3->{red}{black}] (dot_vdd) -- ++ (0,-0.5) node[circ] (p_dot_dc) {} -- ++ (-1,0) -- ++ (0,-0.5) node[pfet, anchor=drain, yscale=-1] (pmos_d) {};
                \draw (p_dot_dc) -- ++ (1,0) -- ++ (0,-0.5)  node[pfet, anchor=drain, yscale=-1] (pmos_a2) {};
                \draw[alt=<3->{red}{black}] (pmos_d.source) -- ++ (0,-0.5) -- ++ (1,0) node[circ] (p_dot_dc_2) {};
                \draw (pmos_a2.source) |- (p_dot_dc_2);  
                
                \draw[alt=<3->{red}{black}] (p_dot_dc_2) -- ++ (0,-0.5) node[circ] (cmos_2) {};
                \draw[alt=<3->{red}{black}] (cmos_1) node[circ] {} -- (cmos_2) {};
                \draw (cmos_2) -- ++ (0,-0.5)  node[nfet, anchor=drain] (nmos_d) {};
                \draw (nmos_d.source) node[nfet, anchor=drain] (nmos_a2) {};
                        
                
                \draw[alt=<3->{red}{black}] (n_dot_c) -- ++ (3,0) node[circ] (dot_ground) {};
                \draw (nmos_a2.source) -- (dot_ground) {};
                
                % Rechte Seite (Inverter)
                \draw (dot_vdd) -- ++ (3,0) -- ++ (0,-1) node[pfet, anchor=drain, yscale=-1] (pmos_inv) {};
                \draw (pmos_inv.source) -- ++ (0,-1) node[circ] (cmos_inv) {};
                \draw (cmos_inv) -- ++ (0,-1) node[nfet, anchor=drain] (nmos_inv) {};
                \draw (nmos_inv.source) -- ++ (0,-1.075) -- (dot_ground) {};
                
                \draw (cmos_2) -- ++ (2.025,0) node[circ] (cmos_3) {};
                \draw (cmos_3) -- (pmos_inv.gate) {};
                \draw (cmos_3) -- (nmos_inv.gate) {};
                
                % Ausgang
                \draw (cmos_inv) -- ++ (1,0) node[ocirc] (g_out) {}
                    node[anchor=west] {\(g(d,c,b,a)\)};
                
                % Versorgungsspannung
                \draw[alt=<3->{red}{black}] (dot_ground) -- ++ (0,-0.5) node[ground] {} node[anchor=east] {\(GND\)};
                
                % Beschriftungen
                \draw (pmos_a.gate) node[anchor=east] {\(a\)};
                \draw (pmos_a2.gate) node[anchor=east] {\(a\)};
                \draw (nmos_a.gate) node[anchor=east] {\(a\)};
                \draw (nmos_a2.gate) node[anchor=east] {\(a\)};
                
                \draw (pmos_b.gate) node[anchor=east] {\(b\)};
                \draw (nmos_b.gate) node[anchor=east] {\(b\)};
                
                \draw (pmos_c.gate) node[anchor=east] {\(c\)};
                \draw (nmos_c.gate) node[anchor=east] {\(\alt<3->{1}{c}\)};
                
                \draw (pmos_d.gate) node[anchor=east] {\(\alt<3->{0}{d}\)};
                \draw (nmos_d.gate) node[anchor=east] {\(d\)};
                
            \end{circuitikz}
        }
        \onslide<3->
        \centering\textcolor{red}{Kurzschluss!}
        }
        \only<5->{
        \begin{solve}
            \resizebox{0.87\textwidth}{!}{
                \begin{circuitikz}[european]
                    % Versorgungsspannung
                    \draw (0,0) node[ocirc] {} node[anchor=east] {\(V_{dd}\)} -- ++ (0,-1) node[circ] (dot_vdd) {};
                    
                    %linke Seite
                    \draw (dot_vdd) -- ++ (-4,0) -- ++ (0,-0.5) node[circ] (p_dot_ab) {};
                    \draw (p_dot_ab) -- ++ (-1,0) node[pfet, anchor=drain, yscale=-1] (pmos_a) {};
                    \draw (p_dot_ab) -- ++ (1,0) node[pfet, anchor=drain, yscale=-1] (pmos_b) {};
                    \draw (pmos_a.source) -- ++ (1,0) node[circ] (p_dot_ab_2) {};
                    \draw (pmos_b.source) -- (p_dot_ab_2);
                    
                    \draw (p_dot_ab_2) -- ++ (0,0) node[pfet, anchor=drain, yscale=-1] (pmos_c) {};
                    \draw (pmos_c.source) -- ++ (0,0) node[circ] (p_c) {}; % Knoten C
                    
            
                    \draw (p_c) -- ++ (0,-0.5) node[circ] (p_dot_dc) {} -- ++ (-1,0) -- ++ (0,-0.5) node[pfet, anchor=drain, yscale=-1] (pmos_d) {};
                    \draw (p_dot_dc) -- ++ (1,0) -- ++ (0,-0.5)  node[pfet, anchor=drain, yscale=-1] (pmos_a2) {};
                    \draw (pmos_d.source) -- ++ (0,-0.5) -- ++ (1,0) node[circ] (p_dot_dc_2) {};
                    \draw (pmos_a2.source) |- (p_dot_dc_2);  
                    \draw (p_dot_dc_2) -- ++ (0,-0.5) node[circ] (cmos_1) {};
                    
                    %Hier anknüpfen
                 
                    
                    \draw (cmos_1) -- ++ (0,-0.5) node[circ] (n_dot_ab) {};
                    \draw (n_dot_ab) -- ++ (-1,0) node[nfet, anchor=drain] (nmos_a) {};
                    \draw (nmos_a.source) -- ++ (0,0) node[nfet, anchor=drain] (nmos_b) {};
                    
                    \draw (n_dot_ab) -- ++ (1,0) -- ++ (0,-0.75) node[nfet, anchor=drain] (nmos_c) {};
                    \draw (nmos_c.source) -- ++ (0,-0.78) node[circ] (n_dot_c) {};
                    \draw (nmos_b.source) -- (n_dot_c)  {};
                    
                    % Mitte
                    
                    
                    
                    \draw (0, -8.1) node[circ] (cmos_2) {};
                    
                    \draw (cmos_1) -- (cmos_2) {};
                    \draw (cmos_2) -- ++ (0,-0.5)  node[nfet, anchor=drain] (nmos_d) {};
                    \draw (nmos_d.source) node[nfet, anchor=drain] (nmos_a2) {};
                            
                    
                    \draw (n_dot_c) -- ++ (3,0) node[circ] (dot_ground) {};
                    \draw (nmos_a2.source) -- (dot_ground) {};
                    
                    % Rechte Seite (Inverter)
                    \draw (dot_vdd) -- ++ (3,0) -- ++ (0,-3) node[pfet, anchor=drain, yscale=-1] (pmos_inv) {};
                    \draw (pmos_inv.source) -- ++ (0,-2.5) node[circ] (cmos_inv) {};
                    \draw (cmos_inv) -- ++ (0,-1) node[nfet, anchor=drain] (nmos_inv) {};
                    \draw (nmos_inv.source) -- ++ (0,-1.075) -- (dot_ground) {};
                    
                    \draw (cmos_2) -- ++ (2.025,0) node[circ] (cmos_3) {};
                    \draw (cmos_3) -- (pmos_inv.gate) {};
                    \draw (cmos_3) -- (nmos_inv.gate) {};
                    
                    % Ausgang
                    \draw (cmos_inv) -- ++ (1,0) node[ocirc] (g_out) {}
                        node[anchor=west] {\(g(d,c,b,a)\)};
                    
                    % Versorgungsspannung
                    \draw (dot_ground) -- ++ (0,-0.5) node[ground] {} node[anchor=east] {\(GND\)};
                    
                    % Beschriftungen
                    \draw (pmos_a.gate) node[anchor=east] {\(a\)};
                    \draw (pmos_a2.gate) node[anchor=east] {\(a\)};
                    \draw (nmos_a.gate) node[anchor=east] {\(a\)};
                    \draw (nmos_a2.gate) node[anchor=east] {\(a\)};
                    
                    \draw (pmos_b.gate) node[anchor=east] {\(b\)};
                    \draw (nmos_b.gate) node[anchor=east] {\(b\)};
                    
                    \draw (pmos_c.gate) node[anchor=east] {\(c\)};
                    \draw (nmos_c.gate) node[anchor=east] {\(c\)};
                    
                    \draw (pmos_d.gate) node[anchor=east] {\(d\)};
                    \draw (nmos_d.gate) node[anchor=east] {\(d\)};
                    
                \end{circuitikz}
            }
        \end{solve}
        }
    \end{column}
    \end{columns}
\end{frame}

\subsection{Aufgabe 2}
\begin{frame}{Wow diese Folien}\onslide<+->
    \begin{columns}[T]
    \begin{column}{0.5\textwidth}
        \begin{exercise}
            Ergänzen Sie\dots
            \only<+->{\\Welche Schaltfunktion?}
        \end{exercise}
        \only<+->{
        \begin{solve}
            \(y = \overline{\ x_5 \ ((\ x_0 \ x_1 \ x_2 \ ) \ \lor \ x_3 \ \lor \ x_4 \ )}\)
            \(\hphantom{y} = \overline{x}_5 \ \lor \ \overline{x}_0 \ \overline{x}_3 \ \overline{x}_4 \ \lor \ \overline{x}_1 \ \overline{x}_3 \ \overline{x}_4 \ \lor \ \overline{x}_2 \ \overline{x}_3 \ \overline{x}_4\)
        \end{solve}
        }
    \end{column}
    \begin{column}{0.5\textwidth}
    \alt<1>{\resizebox{!}{0.8\textheight}{
        \begin{circuitikz}
            \draw
            
            % Versorgungsspannung
            (0,0) node[ocirc] {} node[anchor=east] {\(V_{dd}\)} -- ++ (0,-1) node[circ] (dot1) {}
            %(0,0) node[ocirc] {} node[anchor=east] {$V_{dd}$} -- ++(0,-1) node[circ, color=blue] (dot1) {}
            
            (dot1) -- ++ (-2.5,0) node[circ] (dotx1) {}
            
            % PMOS
            (dot1) -- ++ (2.5, 0)  node[pfet, anchor=drain, yscale=-1] (px5) {}
            (dot1) -- ++ (-5, 0)  node[pfet, anchor=drain, yscale=-1] (px0) {}
            (dotx1) -- ++ (0.0, 0)  node[pfet, anchor=drain, yscale=-1] (px1) {}
            
            (px0.source) -- ++ (1.5, 0) node[pfet, anchor=drain, yscale=-1] (px3) {}
            (px1.source) -- ++ (-1.0, 0) node[circ] {}
            (px3.source) -- ++ (0, 0) node[pfet, anchor=drain, yscale=-1] (px4) {}
            
            (px4.source) -- ++ (0, 0) node[anchor=east] (pOut) {}
            (px5.source) -- ++ (0, -3) node[circ,anchor=north] {}
    
            % Ausgang
            (pOut) -- ++ (1, 0) node[circ, anchor=west] (npY) {} 
            
            (npY) -- ++ (7, 0) node[anchor=west] {\(y\)}
            
            
            % NMOS
            (npY) -- ++ (0,0) node[nfet, anchor=drain] (nx5) {}
            (nx5.source) -- ++ (0,0) node[anchor=west] (nPP) {}
            
            (nPP) -- ++ (-1.5,0) node[nfet, anchor=drain] (nx0) {}
            (nx0.source) -- ++ (0,0) node[nfet, anchor=drain] (nx1) {}
            (nx1.source) -- ++ (0,0) node[nfet, anchor=drain] (nx2) {}
            
            (nx2.source) -- ++ (3.5, 0) node[anchor=south] (nG) {} 
            % GROUND
            (nG) -- ++ (0,-0.5) node[ground] {} node[anchor=east] {\(GND\)}
            
            % Label
            (px0.gate) node[anchor=east] {\(x_0\)}
            (px1.gate) node[anchor=east] {\(x_1\)}
            (px3.gate) node[anchor=east] {\(x_3\)}
            (px4.gate) node[anchor=east] {\(x_4\)}
            (px5.gate) node[anchor=east] {\(x_5\)}
            
            (nx0.gate) node[anchor=east] {\(x_0\)}
            (nx1.gate) node[anchor=east] {\(x_1\)}
            (nx2.gate) node[anchor=east] {\(x_2\)}
            %(nx3.gate) node[anchor=east] {$x_3$}
            %(nx4.gate) node[anchor=east] {$x_4$}
            (nx5.gate) node[anchor=east] {\(x_5\)}
            
            ;
        \end{circuitikz}
    }}{\resizebox{!}{0.8\textheight}{
        \begin{circuitikz}
            \draw
            
            % Versorgungsspannung
            (0,0) node[ocirc] {} node[anchor=east] {\(V_{dd}\)} -- ++ (0,-1) node[circ, color=blue] (dot1) {}
            
            (dot1) -- ++ (-2.5,0) node[circ] (dotx1) {}
            (dot1) -- ++ (2.5, 0)  node[pfet, anchor=drain, yscale=-1] (px5) {}
            
            % PMOS
            ;
            \draw[blue]
            (dot1) -- ++ (0, 0)  node[pfet, anchor=drain, yscale=-1,color=blue] (px2) {}
            
            ;
            \draw
            (dot1) -- ++ (-5, 0)  node[pfet, anchor=drain, yscale=-1] (px0) {}
            (dotx1) -- ++ (0.0, 0)  node[pfet, anchor=drain, yscale=-1] (px1) {}
            (px0.source) -- ++ (1.5, 0) node[pfet, anchor=drain, yscale=-1] (px3) {}
            (px1.source) -- ++ (-1.0, 0) node[circ] {}
            (px5.source) -- ++ (0, -3) node[circ,anchor=north] {}
            ;
            \draw[blue]
            (px2.source) -- ++ (-2.5, 0) node[circ, color=blue] {}
            
            ;
            \draw
            (px3.source) -- ++ (0, 0) node[pfet, anchor=drain, yscale=-1] (px4) {}
            
            (px4.source) -- ++ (0, 0) node[anchor=east] (pOut) {}
            
            % Ausgang
            (pOut) -- ++ (1, 0) node[circ, anchor=west] (npY) {} 
            
            (npY) -- ++ (7, 0) node[anchor=west] {\(y\)}
            
            
            % NMOS
            (npY) -- ++ (0,0) node[nfet, anchor=drain] (nx5) {}
            (nx5.source) -- ++ (0,0) node[circ,color=blue] (nPP) {}
            
            (nPP) -- ++ (-1.5,0) node[nfet, anchor=drain] (nx0) {}
            (nx0.source) -- ++ (0,0) node[nfet, anchor=drain] (nx1) {}
            (nx1.source) -- ++ (0,0) node[nfet, anchor=drain] (nx2) {}
            
            ;
            \draw[blue]
            (nPP) -- ++ (0.5, 0) node[circ, color=blue] (nPP2) {}
            (nPP2) -- ++ (0,0) node[nfet, anchor=drain, color=blue] (nx3) {}
            (nPP2) -- ++ (2.5,0) node[nfet, anchor=drain, color=blue] (nx4) {}		
            
            (nx3.source) -- ++ (0,-3) node[circ, anchor=north, color=blue] {}
            (nx4.source) -- ++ (0,-3) node[anchor=north, color=blue] {}
            
            ;		
            \draw
            (nx2.source) -- ++ (4.5, 0) node[circ, color=blue] (nG) {} 
            % GROUND
            (nG) -- ++ (0,-0.5) node[ground] {} node[anchor=east] {\(GND\)}
            
            % Label
            (px0.gate) node[anchor=east] {\(x_0\)}
            (px1.gate) node[anchor=east] {\(x_1\)}
            (px2.gate) node[anchor=east, color=blue] {\(x_2\)}
            (px3.gate) node[anchor=east] {\(x_3\)}
            (px4.gate) node[anchor=east] {\(x_4\)}
            (px5.gate) node[anchor=east, color=blue] {\(x_5\)}
            
            (nx0.gate) node[anchor=east] {\(x_0\)}
            (nx1.gate) node[anchor=east] {\(x_1\)}
            (nx2.gate) node[anchor=east] {\(x_2\)}
            (nx3.gate) node[anchor=east, color=blue] {\(x_3\)}
            (nx4.gate) node[anchor=east, color=blue] {\(x_4\)}
            (nx5.gate) node[anchor=east] {\(x_5\)}
            
            ;
        \end{circuitikz}
    }}
    \end{column}
    \end{columns}
\end{frame}

\subsection{Aufgabe 3}
\begin{frame}{Was diese?}\onslide<+->
    \begin{exercise}[\only<2->{\(y = a \not\leftrightarrow b \not\leftrightarrow c\)}]
        \resizebox{!}{0.7\textheight}{
            \begin{circuitikz}
                \draw
                
                % Versorgungsspannung
                (0,0) node[ocirc] {} node[anchor=east] {\(V_{dd}\)}
                -- ++ (0,-0.5) node[circ] (dot1) {}
                
                % p-Netz
                (dot1) -- ++ (-1,0) node[pfet, anchor=drain, yscale=-1] (pA) {}
                (pA.source) node[pfet, anchor=drain, yscale=-1] (pBn) {}
                (dot1) -- ++ (1,0) node[pfet, anchor=drain, yscale=-1] (pAn) {}
                (pAn.source) node[pfet, anchor=drain, yscale=-1] (pB) {}
                
                % Mitte, Ausgang
                (pBn.source) node[circ] (dot4){} -- (pB.source) node[circ] (dot3) {}
                (dot3) --++ (1,0) node[circ] (out1) {}
                
                % n-Netz
                (dot4) node[nfet, anchor=drain] (nA) {}
                (nA.source) node[nfet, anchor=drain] (nB) {}
                (dot3) node[nfet, anchor=drain] (nAn) {}
                (nAn.source) node[nfet, anchor=drain] (nBn) {}
                
                % GND
                (nB.source)--++ (1,0) node[circ] (dot5){} -- (nBn.source)
                (dot5) --++ (0,-0.5) node[ground] {} node[anchor=east] {\(GND\)}
                
                % Labels
                (pA.gate) node[anchor=east] {\(a\)}
                (pB.gate) node[anchor=east] {\(b\)}
                (pAn.gate) node[anchor=east] {\(\overline{a}\)}
                (pBn.gate) node[anchor=east] {\(\overline{b}\)}
                
                (nA.gate) node[anchor=east] {\(a\)}
                (nB.gate) node[anchor=east] {\(b\)}
                (nAn.gate) node[anchor=east] {\(\overline{a}\)}
                (nBn.gate) node[anchor=east] {\(\overline{b}\)}
                
                
                
                
                
                (out1) --++ (0.5,0) node[circ] (dot7){}
                ++ (1,0) node[circ] (dot6){}
                --++ (1,0) node[circ] (out2){}
                
                % p-Netz
                (dot6) node[pfet, anchor=source,yscale=-1] (pA){}
                (pA.gate) -| (dot7)
                % n-Netz
                (dot6) node[nfet, anchor=drain] (nA){}
                (nA.gate) -| (dot7)
                
                (pA.drain) --++ (0,0.5) node[ocirc] {} node[anchor=east] {\(V_{dd}\)}
                (nA.source) --++ (0,-0.5) node[ground] {} node[anchor=east] {\(GND\)}
                
                
                
                
                % Mitte
                (out2) ++ (2,0) node[circ] (dot8){}
                -- ++ (2,0) node[circ] (dot9){}
                (dot9) -- ++ (1,0) node[ocirc] (out3){} node[anchor=west] {\(y\)}
                
                
                % p-Netz
                (dot8) node[pfet, anchor=source, yscale=-1] (pCn) {}
                (pCn.drain) node[pfet, anchor=source, yscale=-1] (pO) {}
                (out1) -- ++ (0,3) -| (pO.gate)
                
                (dot9) node[pfet, anchor=source, yscale=-1] (pOn) {}
                (pOn.drain) node[pfet, anchor=source, yscale=-1] (pC) {}
                (out2) -- ++ (0,0.3) -| (pOn.gate)
                
                
                % oben
                (pO.drain) -- ++ (1,0) node[circ] (dot10) {} -- (pC.drain)
                (dot10) -- ++ (0,0.5) node[ocirc] {} node[anchor=east] {\(V_{dd}\)}
                
                
                % n-Netz
                (dot8) node[nfet, anchor=drain] (nO) {}
                (nO.source) node[nfet, anchor=drain] (nCn) {}
                (out2) |- (nO.gate)
                
                (dot9) node[nfet, anchor=drain] (nC) {}
                (nC.source) node[nfet, anchor=drain] (nOn) {}
                (out1) -- ++ (0,-3) -| (nOn.gate)
                
                % unten
                (nCn.source) |- ++ (1,-0.3) node[circ] (dot11) {} -| (nOn.source)
                (dot11) -- ++ (0,-0.5) node[ground] {} node[anchor=east] {\(GND\)}
                
                
                % Labels
                (pCn.gate) node[anchor=east] {\(\overline{c} \)}
                (pC.gate) node[anchor=east] {\(c\)}
                (nCn.gate) node[anchor=east] {\(\overline{c} \)}
                (nC.gate) node[anchor=east] {\(c\)}
                
                
                ;
        \end{circuitikz}
        }
    \end{exercise}
\end{frame}

\end{document}% chktex 17